\documentclass[12pt, a4paper]{article}

% --- Pakete ---
\usepackage[ngerman]{babel}
\usepackage[utf8]{inputenc}
\usepackage[T1]{fontenc}
\usepackage{lmodern}
\usepackage{geometry}
\usepackage{graphicx}
\usepackage{hyperref}
\usepackage{listings}
\usepackage{xcolor}
\usepackage{booktabs}
\usepackage{enumitem}
\usepackage{parskip}
\usepackage{fancyhdr}

\geometry{left=2.5cm, right=2.5cm, top=2.5cm, bottom=2.5cm}

% --- Farben ---
\definecolor{codeblue}{RGB}{0,102,204}
\definecolor{codegray}{RGB}{100,100,100}
\definecolor{codegreen}{RGB}{0,128,0}
\definecolor{backgray}{RGB}{245,245,245}

% --- Listings-Stil ---
\lstset{
    backgroundcolor=\color{backgray},
    basicstyle=\ttfamily\small,
    keywordstyle=\color{codeblue}\bfseries,
    commentstyle=\color{codegreen},
    stringstyle=\color{codegray},
    numbers=left,
    numberstyle=\tiny\color{codegray},
    numbersep=8pt,
    frame=single,
    breaklines=true,
    showstringspaces=false,
    tabsize=4,
    captionpos=b,
}

% --- Header/Footer ---
\pagestyle{fancy}
\fancyhf{}
\fancyhead[L]{\textbf{KuchenApp -- Projektdokumentation}}
\fancyhead[R]{\thepage}
\renewcommand{\headrulewidth}{0.4pt}

% --- Hyperlinks ---
\hypersetup{
    colorlinks=true,
    linkcolor=codeblue,
    urlcolor=codeblue,
}

% =====================================================================
\begin{document}

% --- Titelseite ---
\begin{titlepage}
\centering
\vspace*{3cm}
{\Huge\textbf{KuchenApp}\\[0.5cm]}
{\LARGE Projektdokumentation}\\[2cm]
{\Large Desktop-Anwendung zur Verwaltung von\\Kuchen-Listen und IServ-Mail-Versand}\\[3cm]
{\large Erstellt mit PySide6 (Qt for Python)}\\[1cm]
{\large Python 3.10+}\\[3cm]
{\large\today}
\vfill
\end{titlepage}

\tableofcontents
\newpage

% =====================================================================
\section{Projektübersicht}

\textbf{KuchenApp} ist eine Desktop-Anwendung, die im Rahmen eines Schulprojekts
(Berufsschule, Fachinformatiker Anwendungsentwicklung) entwickelt wurde. Die Anwendung
dient zur Verwaltung einer Kuchen-Liste für eine Schulklasse und ermöglicht den
automatisierten Versand von Erinnerungs-E-Mails über das IServ-System.

\subsection{Funktionsumfang}
\begin{itemize}
    \item CSV-Dateien importieren, bearbeiten und speichern
    \item Dynamische Tabellendarstellung mit Sortier- und Suchfunktion
    \item Aufsteigende und absteigende Sortierung per Radio-Buttons
    \item Verwaltung einer Empfänger-Mailliste (separater Dialog)
    \item Automatische HTML-Tabellengenerierung für E-Mail-Inhalte
    \item E-Mail-Versand über IServ-SMTP mit Login-Dialog
    \item INI-basierte Einstellungsverwaltung
    \item Standalone-Build für Linux und Windows (PyInstaller)
    \item CI/CD über GitHub Actions
\end{itemize}

\subsection{Technologie-Stack}
\begin{table}[h!]
\centering
\begin{tabular}{ll}
\toprule
\textbf{Komponente} & \textbf{Technologie} \\
\midrule
Programmiersprache & Python 3.10+ \\
GUI-Framework & PySide6 (Qt 6 for Python) \\
UI-Design & Qt Designer (.ui-Dateien) \\
Datenpersistenz & CSV (comma-separated values) \\
Einstellungen & INI-Dateien (configparser) \\
E-Mail-Versand & smtplib (SMTP\_SSL) \\
Build/Packaging & PyInstaller \\
CI/CD & GitHub Actions \\
\bottomrule
\end{tabular}
\caption{Verwendete Technologien}
\end{table}


% =====================================================================
\section{Projektstruktur}

Das Projekt folgt einer klaren Trennung zwischen Entwicklungs- und Release-Code.

\subsection{Verzeichnisstruktur}

\begin{lstlisting}[language={}, caption={Gesamtstruktur}]
Kuchen/
|-- build_release.py          # Build-Skript
|-- requirements.txt          # Python-Abhängigkeiten
|-- .github/workflows/        # CI/CD Pipeline
|   +-- build_release.yml
|-- develop/                  # Entwicklungsumgebung
|   |-- KuchenMain.py         # Hauptprogramm (Entry Point)
|   |-- src/                  # Quellcode-Module
|   |   |-- __init__.py
|   |   |-- AppLogger.py
|   |   |-- DataManager.py
|   |   |-- IServAPIEdited_standalone.py
|   |   |-- IservMailManager.py
|   |   |-- KuchenMailManager.py
|   |   |-- LoginDialog.py
|   |   +-- SettingsManager.py
|   |-- ui/                   # UI-Dateien
|   |   |-- __init__.py
|   |   |-- UIMainWindow.ui / _ui.py
|   |   |-- UIMailsDialog.ui / _ui.py
|   |   +-- UILoginDialog.ui / _ui.py
|   |-- assets/
|   |   +-- kuchen_icon.svg
|   +-- Data/
|       |-- CakeData.csv
|       |-- StandardCakeData.csv
|       |-- Settings/settings.ini
|       +-- Klassen/EITB23A.csv
+-- release/                  # Generierte Releases
    |-- KuchenApp/            # Source-Bundle
    |-- KuchenApp.zip
    |-- KuchenApp_standalone_linux/
    |-- KuchenApp_standalone_linux.zip
    |-- KuchenApp_standalone_windows/
    +-- KuchenApp_standalone_windows.zip
\end{lstlisting}

\subsection{Trennung: develop vs. release}

\begin{description}
    \item[develop/] Enthält den vollständigen Quellcode, .ui-Designdateien
        und Testdaten. Hier wird aktiv entwickelt.
    \item[release/] Wird durch \texttt{build\_release.py} automatisch generiert.
        Enthält nur die für den Endbenutzer nötigen Dateien. Die CSV-Dateien
        werden mit leeren Standard-Headern ausgeliefert.
\end{description}


% =====================================================================
\section{Architektur und Entwurfsmuster}

\subsection{Gesamtarchitektur}

Die Anwendung folgt einer \textbf{Schichtenarchitektur}:

\begin{enumerate}
    \item \textbf{Präsentationsschicht} -- GUI-Klassen (\texttt{CMainWindow},
        \texttt{CKuchenDialog}, \texttt{CLoginDialog}), Qt-UI-Dateien
    \item \textbf{Logikschicht} -- Datenverarbeitung (\texttt{DataManager}),
        Einstellungen (\texttt{SettingsManager}), Mail-Versand
        (\texttt{IservMailManager})
    \item \textbf{Datenzugriffsschicht} -- CSV-Dateien, INI-Konfiguration,
        IServ-SMTP-API
\end{enumerate}

\subsection{Verwendete Entwurfsmuster}

\subsubsection{Singleton-Pattern (AppLogger)}

Der \texttt{AppLogger} implementiert das \textbf{Singleton-Pattern} über Pythons
\texttt{\_\_new\_\_}-Methode. Damit existiert während der gesamten Laufzeit
exakt eine Logger-Instanz.

\textbf{Fachliche Begründung:} Alle Klassen sollen Fehlermeldungen an die GUI
weiterleiten können, ohne eine Referenz auf das Hauptfenster zu benötigen. Der
Singleton stellt sicher, dass alle \texttt{AppLogger()}-Aufrufe dieselbe Instanz
zurückgeben und somit die gleichen Qt-Signale nutzen.

\begin{lstlisting}[language=Python, caption={Singleton-Implementierung}]
class AppLogger(QObject):
    errorOccurred = Signal(str, str)
    warningOccurred = Signal(str, str)
    infoOccurred = Signal(str, str)

    _instance = None

    def __new__(cls):
        if cls._instance is None:
            cls._instance = super().__new__(cls)
        return cls._instance

    def __init__(self):
        if not hasattr(self, '_initialized'):
            super().__init__()
            self._initialized = True
\end{lstlisting}

\subsubsection{Observer-Pattern (Qt Signal/Slot)}

Die Kommunikation zwischen den Schichten erfolgt über das
\textbf{Observer-Pattern}, implementiert durch Qt's Signal/Slot-Mechanismus.

\begin{itemize}
    \item \texttt{AppLogger} emittiert Signale (\texttt{errorOccurred},
        \texttt{warningOccurred}, \texttt{infoOccurred})
    \item \texttt{CMainWindow} verbindet diese mit Slot-Methoden, die
        \texttt{QMessageBox}-Dialoge anzeigen
    \item \texttt{QStandardItemModel.dataChanged} benachrichtigt bei
        Zellenänderungen in der Tabelle
\end{itemize}

\textbf{Fachliche Begründung:} Lose Kopplung -- der \texttt{DataManager} und
\texttt{IservMailManager} kennen die GUI nicht, sondern kommunizieren über den
\texttt{AppLogger}. Dies ermöglicht eine saubere Trennung der Schichten.

\subsubsection{Model-View-Pattern (Qt MVC)}

Die Tabellendarstellung nutzt Qt's \textbf{Model-View}-Architektur:

\begin{itemize}
    \item \textbf{Model:} \texttt{QStandardItemModel} -- hält die Tabelleneinträge
    \item \textbf{View:} \texttt{QTableView} -- zeigt die Daten an
    \item \textbf{Controller:} Die jeweilige Window-Klasse vermittelt zwischen
        \texttt{DataManager} und dem Model
\end{itemize}


% =====================================================================
\section{Module im Detail}

\subsection{KuchenMain.py -- Hauptfenster}

\texttt{CMainWindow} erbt von \texttt{QMainWindow} und \texttt{Ui\_MainWindow}
(Multiple Inheritance, empfohlener Qt-Ansatz). Es ist der \textbf{Entry Point}
der Anwendung.

\textbf{Schlüsselfunktionen:}

\begin{description}[style=nextline]
    \item[\texttt{\_createRadioButtons()}]
        Erzeugt dynamisch Radio-Buttons für jede CSV-Spalte (Sortierung)
        sowie ein Ascending/Descending-Paar. Verwendet \texttt{QButtonGroup}
        für die Sortierrichtung.
    \item[\texttt{SortColumn()}]
        Liest den ausgewählten Radio-Button aus, bestimmt den Spaltenindex
        und ruft \texttt{DataManager.sortData()} mit dem
        \texttt{aReverse}-Parameter auf. Berücksichtigt aktive Suchfilter.
    \item[\texttt{Search(aText)}]
        Filtert die Originaldaten (\texttt{mData}) nach dem Suchbegriff.
        Ist ein Sortier-Radio-Button aktiv, wird nur in der entsprechenden
        Spalte gesucht. Ansonsten: Volltextsuche über alle Spalten.
        Nach dem Filtern wird die Sortierung erneut angewandt.
    \item[\texttt{CellEdited(aItem)}]
        Synchronisiert Änderungen in der Tabelle zurück in \texttt{mSortedData}
        \emph{und} \texttt{mData}. Sucht die Original-Zeile anhand aller
        anderen Spaltenwerte (da die Reihenfolge abweichen kann).
    \item[\texttt{FillTable()}]
        Löscht das \texttt{QStandardItemModel}, setzt Header und füllt die
        Tabelle aus \texttt{mSortedData}. Wird nach jeder Datenänderung
        aufgerufen.
    \item[\texttt{SendMails()}]
        Zeigt einen Bestätigungsdialog mit der Empfängeranzahl, öffnet den
        Login-Dialog und wertet den Tri-State-Rückgabewert aus
        (\texttt{None}=abgebrochen, \texttt{True}=gesendet,
        \texttt{False}=fehlgeschlagen).
    \item[\texttt{ImportFile()}]
        Öffnet einen \texttt{QFileDialog}, importiert die CSV und
        aktualisiert die Radio-Buttons dynamisch anhand der neuen Header.
\end{description}

\textbf{PyInstaller-Kompatibilität:} Im \texttt{\_\_main\_\_}-Block wird
\texttt{sys.frozen} geprüft, um den Icon-Pfad korrekt aufzulösen
(\texttt{sys.\_MEIPASS} im Bundle vs. \texttt{\_\_file\_\_} im
Entwicklungsmodus).


\subsection{DataManager.py -- Datenverwaltung}

Der \texttt{DataManager} ist die zentrale Klasse für alle Datenoperationen.
Er verwaltet zwei unabhängige Datensätze:

\begin{itemize}
    \item \textbf{Haupttabelle} (Kuchen-Daten): \texttt{mData}, \texttt{mSortedData},
        \texttt{mMainHeaders}
    \item \textbf{Mail-Tabelle} (Empfänger): \texttt{mMailData}, \texttt{mSortedMailData},
        \texttt{mMailHeaders}
\end{itemize}

\textbf{Duale Datenhaltung:} \texttt{mData} enthält immer die \emph{unsortierte}
Originaldaten. \texttt{mSortedData} ist eine gefilterte/sortierte Kopie, die
in der Tabelle angezeigt wird. Dies ermöglicht das Zurücksetzen von
Sortierung und Suche.

\textbf{Schlüsselfunktionen:}

\begin{description}[style=nextline]
    \item[\texttt{LoadStandardFile(aForMainTable)}]
        Lädt die CSV-Datei, deren Name in den Settings steht. Verwendet
        \texttt{csv.Sniffer} zur automatischen Dialekt-Erkennung (Trennzeichen,
        Quoting). Extrahiert die erste Zeile als Header.
    \item[\texttt{ImportFile(aPath)}]
        Importiert eine externe CSV. Speichert den Dateinamen in
        \texttt{mImportedFilename}. Beim nächsten \texttt{SaveFile()} wird
        dieser Name in die Settings übernommen, sodass die importierte Datei
        beim nächsten Start automatisch geladen wird.
    \item[\texttt{SaveFile(aForMainTable)}]
        Speichert die Daten als CSV. Wurde zuvor eine Datei importiert,
        aktualisiert \texttt{SaveFile} automatisch die Settings mit dem
        neuen Dateinamen.
    \item[\texttt{sortData(aSortColumnIndex, aReverse, aForMain)}]
        Sortiert \texttt{mSortedData} (oder \texttt{mSortedMailData}) nach
        der gegebenen Spalte. Der Parameter \texttt{aReverse} steuert die
        Richtung (aufsteigend/absteigend).
    \item[\texttt{\_padRows(aData, aHeaderLen)}]
        Füllt Datenzeilen mit leeren Strings auf, falls sie weniger Spalten
        als der Header haben. Dies erhöht die Robustheit beim Import
        unvollständiger CSVs.
    \item[\texttt{createCompleteMailData()}]
        Bereitet die Mail-Daten auf: extrahiert E-Mail-Adressen und generiert
        einen HTML-String der Kuchen-Tabelle für den E-Mail-Body.
    \item[\texttt{createHtmlDataString()}]
        Erzeugt eine formatierte HTML-Tabelle mit Inline-CSS. Verwendet
        \texttt{html.escape()} zur Vermeidung von XSS-Schwachstellen.
\end{description}

\textbf{PyInstaller-Kompatibilität:} Die Funktion \texttt{\_get\_base\_dir()}
gibt bei \texttt{sys.frozen == True} das Verzeichnis der ausführbaren Datei
zurück, ansonsten das Projektverzeichnis.


\subsection{SettingsManager.py -- Einstellungsverwaltung}

Verwaltet die Anwendungseinstellungen in einer INI-Datei
(\texttt{Data/Settings/settings.ini}) über Pythons \texttt{configparser}-Modul.

\textbf{Konfigurationsstruktur:}

\begin{lstlisting}[language={}, caption={settings.ini}]
[Klassen]
Dateiname = StandardKlasse.csv

[CakeData]
Dateiname = StandardCakeData.csv

[IServ]
Domain = wvss.de
\end{lstlisting}

\textbf{Default-Werte:} Die Klasse definiert ein \texttt{\_DEFAULTS}-Dictionary.
Beim Instanziieren werden zuerst die Defaults geladen, dann die INI-Datei
darübergelegt. So sind immer Fallback-Werte vorhanden.


\subsection{AppLogger.py -- Singleton-Logger}

Implementiert das Singleton-Pattern (siehe Abschnitt 3.2.1). Stellt drei
Schweregrade bereit:

\begin{itemize}
    \item \texttt{error(aTitle, aMessage)} -- Kritische Fehler (roter Dialog)
    \item \texttt{warning(aTitle, aMessage)} -- Warnungen (gelber Dialog)
    \item \texttt{info(aTitle, aMessage)} -- Informationen (blauer Dialog)
\end{itemize}

Jeder Aufruf emittiert ein Qt-Signal mit Titel und Nachricht. Die GUI
verbindet diese Signale mit \texttt{QMessageBox.critical()},
\texttt{QMessageBox.warning()} bzw. \texttt{QMessageBox.information()}.


\subsection{KuchenMailManager.py -- Mail-Listen-Dialog}

\texttt{CKuchenDialog} ist ein modaler Dialog (\texttt{QDialog}) zur Verwaltung
der Empfänger-Mailliste. Er teilt sich den \texttt{DataManager} mit dem
Hauptfenster und arbeitet auf den Mail-Daten (\texttt{mMailData}).

Die Funktionalität ist analog zum Hauptfenster: dynamische Radio-Buttons,
Sortierung (ascending/descending), Suche, Import, Speichern, Zeilen
hinzufügen/löschen.


\subsection{LoginDialog.py -- IServ-Login}

\texttt{CLoginDialog} ist ein modaler Dialog für die SMTP-Authentifizierung.

\textbf{Tri-State-Logik:}
\begin{itemize}
    \item \texttt{mSendState = None} -- Benutzer hat abgebrochen
    \item \texttt{mSendState = True} -- Mails erfolgreich gesendet
    \item \texttt{mSendState = False} -- Fehler beim Senden (Dialog bleibt offen)
\end{itemize}


\subsection{IservMailManager.py -- Mail-Versand}

Erstellt eine Verbindung zur IServ-API und sendet E-Mails über SMTP\_SSL.
Das Passwort wird im \texttt{finally}-Block auf \texttt{None} gesetzt,
um es nicht unnötig im Speicher zu halten (\textbf{Security Best Practice}).

Fehlerbehandlung fängt spezifische Exceptions:
\texttt{OSError}, \texttt{smtplib.SMTPException}, \texttt{ValueError},
\texttt{ConnectionError}.


\subsection{IServAPIEdited\_standalone.py -- IServ-API}

Eine angepasste Version der IServ-API, die \textbf{ausschließlich} die
Python-Standardbibliothek verwendet (kein \texttt{pip install} nötig).

Ersetzt externe Abhängigkeiten:
\begin{itemize}
    \item \texttt{requests} $\rightarrow$ \texttt{urllib.request}
    \item \texttt{beautifulsoup4} $\rightarrow$ eigene \texttt{\_HTMLNode}-Klasse
    \item \texttt{smtplib} für SMTP\_SSL-Verbindungen
\end{itemize}

Für den E-Mail-Versand wird \texttt{email.mime} zur MIME-Nachrichtenerstellung
und \texttt{smtplib.SMTP\_SSL} für die verschlüsselte Verbindung genutzt.


% =====================================================================
\section{Datenpersistenz}

\subsection{CSV-Format}

Die Anwendung speichert Daten als CSV-Dateien (Comma-Separated Values).
Beim Laden wird \texttt{csv.Sniffer} zur automatischen Erkennung des
Trennzeichens verwendet -- so werden sowohl Komma- als auch
Semikolon-getrennte Dateien korrekt gelesen.

\textbf{StandardCakeData.csv} (leere Vorlage):
\begin{lstlisting}[language={}, caption={CSV-Header}]
Name,CakeCount,Hanuta,Waffel,Date
\end{lstlisting}

\textbf{Dynamische Header:} Die Header werden nicht hartcodiert, sondern
aus der ersten Zeile der CSV-Datei gelesen. Die GUI-Radio-Buttons für
die Sortierung werden dynamisch aus diesen Headern generiert.

\subsection{Settings (INI)}

Einstellungen werden über Pythons \texttt{configparser}-Modul in einer
INI-Datei verwaltet. Die Datei wird mit Defaults initialisiert und bei
Änderungen (z.\,B. Import einer neuen CSV) automatisch aktualisiert.


% =====================================================================
\section{Build-System}

\subsection{build\_release.py}

Das zentrale Build-Skript automatisiert die Release-Erstellung:

\begin{lstlisting}[language=bash, caption={Verwendung}]
python build_release.py             # Alle 3 Versionen
python build_release.py --source    # Nur Source-Bundle
python build_release.py --linux     # Nur Linux Standalone
python build_release.py --windows   # Nur Windows Standalone
\end{lstlisting}

\textbf{Build-Schritte:}
\begin{enumerate}
    \item \textbf{Clean:} Alter \texttt{release/}-Ordner wird gelöscht
    \item \textbf{UI-Kompilierung:} \texttt{.ui} $\rightarrow$ \texttt{\_ui.py}
        via \texttt{pyside6-uic}
    \item \textbf{Kopieren:} Quelldateien werden in die Release-Struktur kopiert
    \item \textbf{README:} Automatische README.txt-Generierung
    \item \textbf{ZIP:} Source-Bundle als ZIP verpackt
    \item \textbf{PyInstaller:} Standalone-Binary für das aktuelle OS
    \item \textbf{Remote-Build:} Für das andere OS via GitHub Actions
        (mit \texttt{gh} CLI)
\end{enumerate}

\subsection{PyInstaller -- Standalone-Builds}

PyInstaller verpackt Python-Interpreter, PySide6 und alle Quelldateien in
eine einzelne ausführbare Datei (\texttt{--onefile}).

\textbf{Wichtige Optionen:}
\begin{itemize}
    \item \texttt{--onefile} -- Einzelne Datei statt Ordner
    \item \texttt{--windowed} -- Kein Konsolenfenster (GUI-App)
    \item \texttt{--add-data} -- Icon und Data-Verzeichnis einbetten
\end{itemize}

\textbf{Laufzeit-Pfadauflösung:} Da PyInstaller Dateien in ein temporäres
Verzeichnis (\texttt{sys.\_MEIPASS}) entpackt, aber \texttt{Data/}
beschreibbar sein muss, wird \texttt{Data/} neben die Executable kopiert.
Die Funktion \texttt{\_get\_base\_dir()} in \texttt{DataManager} und
\texttt{SettingsManager} unterscheidet automatisch:

\begin{lstlisting}[language=Python, caption={Pfadauflösung}]
def _get_base_dir():
    if getattr(sys, 'frozen', False):
        return os.path.dirname(sys.executable)
    return os.path.dirname(os.path.dirname(
        os.path.abspath(__file__)))
\end{lstlisting}


\subsection{GitHub Actions -- CI/CD}

Der Workflow \texttt{.github/workflows/build\_release.yml} baut automatisch
für beide Betriebssysteme:

\begin{itemize}
    \item \textbf{Trigger:} Version-Tag (\texttt{v*}) oder manueller Dispatch
    \item \textbf{Matrix-Build:} Ubuntu + Windows parallel
    \item \textbf{Release-Job:} Erstellt ein GitHub Release mit allen 3 ZIPs
        als Download-Anhänge
\end{itemize}

\begin{lstlisting}[language={}, caption={Release auslösen}]
git tag v1.0
git push --tags
\end{lstlisting}


% =====================================================================
\section{Sicherheitsaspekte}

\begin{itemize}
    \item \textbf{HTML-Escaping:} Alle Tabelleninhalte werden mit
        \texttt{html.escape()} in der Mail-HTML-Generierung escaped,
        um Cross-Site-Scripting (XSS) zu verhindern.
    \item \textbf{Passwort-Löschung:} Das SMTP-Passwort wird im
        \texttt{finally}-Block von \texttt{IservMailManager.sendMail()}
        auf \texttt{None} gesetzt.
    \item \textbf{Spezifische Exceptions:} Statt generischem
        \texttt{except Exception} werden nur erwartete Exception-Typen
        gefangen (\texttt{OSError}, \texttt{csv.Error},
        \texttt{smtplib.SMTPException}, etc.).
    \item \textbf{SMTP\_SSL:} Die Verbindung zu IServ erfolgt über
        SSL-verschlüsseltes SMTP (Port 465).
\end{itemize}


% =====================================================================
\section{Benutzerhandbuch}

\subsection{Installation (Source-Version)}

\begin{lstlisting}[language=bash]
pip install PySide6
cd KuchenApp
python KuchenMain.py
\end{lstlisting}

\subsection{Installation (Standalone)}

Die Standalone-Version benötigt keine Installation:
\begin{itemize}
    \item \textbf{Linux:} \texttt{./KuchenApp}
    \item \textbf{Windows:} Doppelklick auf \texttt{KuchenApp.exe}
\end{itemize}

\subsection{Bedienung}

\begin{enumerate}
    \item \textbf{Daten laden:} Beim Start wird die zuletzt verwendete CSV
        automatisch geladen.
    \item \textbf{Import:} Über den \emph{Import}-Button kann eine andere
        CSV-Datei geladen werden. Die Spalten und Sortier-Buttons passen
        sich automatisch an.
    \item \textbf{Bearbeiten:} Zellen können direkt in der Tabelle
        bearbeitet werden (Doppelklick).
    \item \textbf{Zeilen hinzufügen/löschen:} Über die Buttons
        \emph{Add New} und \emph{Delete}.
    \item \textbf{Sortierung:} Radio-Buttons rechts neben der Tabelle.
        Ascending/Descending wählt die Richtung.
    \item \textbf{Suche:} Das Suchfeld filtert live. Bei aktivem
        Sortier-Button wird nur in der entsprechenden Spalte gesucht.
    \item \textbf{Speichern:} Über den \emph{Save}-Button. Bei importierten
        Dateien wird der neue Dateiname in den Settings gespeichert.
    \item \textbf{Mail-Liste:} Über \emph{View/Edit Mail List} öffnet sich
        ein separater Dialog zur Verwaltung der Empfänger.
    \item \textbf{Mails senden:} Über \emph{SendMails}. Es erscheint eine
        Bestätigung mit der Empfängeranzahl, dann der Login-Dialog.
\end{enumerate}


% =====================================================================
\section{Glossar}

\begin{description}
    \item[CSV] Comma-Separated Values -- Tabellenformat mit Trennzeichen
    \item[INI] Initialisierungsdatei -- Einfaches Konfigurationsformat
    \item[PySide6] Offizielle Python-Bindings für das Qt 6 Framework
    \item[Qt Designer] Visueller Editor für Qt-Oberfl\"achen (.ui-Dateien)
    \item[PyInstaller] Tool zum Verpacken von Python-Programmen als Standalone
    \item[SMTP\_SSL] Simple Mail Transfer Protocol über SSL-Verschlüsselung
    \item[Singleton] Entwurfsmuster: Garantiert genau eine Instanz einer Klasse
    \item[Observer] Entwurfsmuster: Objekte werden über Zustandsänderungen
        benachrichtigt (in Qt: Signal/Slot)
    \item[MVC] Model-View-Controller -- Architekturmuster zur Trennung von
        Daten, Darstellung und Steuerung
    \item[CI/CD] Continuous Integration / Continuous Deployment -- Automatisierte
        Build- und Release-Pipelines
    \item[XSS] Cross-Site Scripting -- Sicherheitslücke durch unescapten
        HTML-Input
    \item[IServ] Schulplattform mit E-Mail, Kalender und Dateiverwaltung
\end{description}

\end{document}
